% Options for packages loaded elsewhere
% Options for packages loaded elsewhere
\PassOptionsToPackage{unicode}{hyperref}
\PassOptionsToPackage{hyphens}{url}
%
\documentclass[
  letterpaper,
  11pt,
  english,
  onehalfspacing,
  headsepline]{MastersDoctoralThesis}
\usepackage{xcolor}
\usepackage{amsmath,amssymb}
\setcounter{secnumdepth}{5}
\usepackage{iftex}
\ifPDFTeX
  \usepackage[T1]{fontenc}
  \usepackage[utf8]{inputenc}
  \usepackage{textcomp} % provide euro and other symbols
\else % if luatex or xetex
  \usepackage{unicode-math} % this also loads fontspec
  \defaultfontfeatures{Scale=MatchLowercase}
  \defaultfontfeatures[\rmfamily]{Ligatures=TeX,Scale=1}
\fi
\usepackage{lmodern}
\ifPDFTeX\else
  % xetex/luatex font selection
\fi
% Use upquote if available, for straight quotes in verbatim environments
\IfFileExists{upquote.sty}{\usepackage{upquote}}{}
\IfFileExists{microtype.sty}{% use microtype if available
  \usepackage[]{microtype}
  \UseMicrotypeSet[protrusion]{basicmath} % disable protrusion for tt fonts
}{}
\makeatletter
\@ifundefined{KOMAClassName}{% if non-KOMA class
  \IfFileExists{parskip.sty}{%
    \usepackage{parskip}
  }{% else
    \setlength{\parindent}{0pt}
    \setlength{\parskip}{6pt plus 2pt minus 1pt}}
}{% if KOMA class
  \KOMAoptions{parskip=half}}
\makeatother
% Make \paragraph and \subparagraph free-standing
\makeatletter
\ifx\paragraph\undefined\else
  \let\oldparagraph\paragraph
  \renewcommand{\paragraph}{
    \@ifstar
      \xxxParagraphStar
      \xxxParagraphNoStar
  }
  \newcommand{\xxxParagraphStar}[1]{\oldparagraph*{#1}\mbox{}}
  \newcommand{\xxxParagraphNoStar}[1]{\oldparagraph{#1}\mbox{}}
\fi
\ifx\subparagraph\undefined\else
  \let\oldsubparagraph\subparagraph
  \renewcommand{\subparagraph}{
    \@ifstar
      \xxxSubParagraphStar
      \xxxSubParagraphNoStar
  }
  \newcommand{\xxxSubParagraphStar}[1]{\oldsubparagraph*{#1}\mbox{}}
  \newcommand{\xxxSubParagraphNoStar}[1]{\oldsubparagraph{#1}\mbox{}}
\fi
\makeatother


\usepackage{longtable,booktabs,array}
\usepackage{calc} % for calculating minipage widths
% Correct order of tables after \paragraph or \subparagraph
\usepackage{etoolbox}
\makeatletter
\patchcmd\longtable{\par}{\if@noskipsec\mbox{}\fi\par}{}{}
\makeatother
% Allow footnotes in longtable head/foot
\IfFileExists{footnotehyper.sty}{\usepackage{footnotehyper}}{\usepackage{footnote}}
\makesavenoteenv{longtable}
\usepackage{graphicx}
\makeatletter
\newsavebox\pandoc@box
\newcommand*\pandocbounded[1]{% scales image to fit in text height/width
  \sbox\pandoc@box{#1}%
  \Gscale@div\@tempa{\textheight}{\dimexpr\ht\pandoc@box+\dp\pandoc@box\relax}%
  \Gscale@div\@tempb{\linewidth}{\wd\pandoc@box}%
  \ifdim\@tempb\p@<\@tempa\p@\let\@tempa\@tempb\fi% select the smaller of both
  \ifdim\@tempa\p@<\p@\scalebox{\@tempa}{\usebox\pandoc@box}%
  \else\usebox{\pandoc@box}%
  \fi%
}
% Set default figure placement to htbp
\def\fps@figure{htbp}
\makeatother


% definitions for citeproc citations
\NewDocumentCommand\citeproctext{}{}
\NewDocumentCommand\citeproc{mm}{%
  \begingroup\def\citeproctext{#2}\cite{#1}\endgroup}
\makeatletter
 % allow citations to break across lines
 \let\@cite@ofmt\@firstofone
 % avoid brackets around text for \cite:
 \def\@biblabel#1{}
 \def\@cite#1#2{{#1\if@tempswa , #2\fi}}
\makeatother
\newlength{\cslhangindent}
\setlength{\cslhangindent}{1.5em}
\newlength{\csllabelwidth}
\setlength{\csllabelwidth}{3em}
\newenvironment{CSLReferences}[2] % #1 hanging-indent, #2 entry-spacing
 {\begin{list}{}{%
  \setlength{\itemindent}{0pt}
  \setlength{\leftmargin}{0pt}
  \setlength{\parsep}{0pt}
  % turn on hanging indent if param 1 is 1
  \ifodd #1
   \setlength{\leftmargin}{\cslhangindent}
   \setlength{\itemindent}{-1\cslhangindent}
  \fi
  % set entry spacing
  \setlength{\itemsep}{#2\baselineskip}}}
 {\end{list}}
\usepackage{calc}
\newcommand{\CSLBlock}[1]{\hfill\break\parbox[t]{\linewidth}{\strut\ignorespaces#1\strut}}
\newcommand{\CSLLeftMargin}[1]{\parbox[t]{\csllabelwidth}{\strut#1\strut}}
\newcommand{\CSLRightInline}[1]{\parbox[t]{\linewidth - \csllabelwidth}{\strut#1\strut}}
\newcommand{\CSLIndent}[1]{\hspace{\cslhangindent}#1}



\setlength{\emergencystretch}{3em} % prevent overfull lines

\providecommand{\tightlist}{%
  \setlength{\itemsep}{0pt}\setlength{\parskip}{0pt}}



 


% -----------------------------------------------------------------------
%  Extra packages and thesis metadata
%  (MastersDoctoralThesis.cls loads: geometry, hyperref, caption, etc.)
% -----------------------------------------------------------------------

\usepackage[utf8]{inputenc}
\usepackage[T1]{fontenc}
\usepackage{mathpazo}          % Palatino font (matches the template)
\usepackage{csquotes}          % Language-aware quotation marks

% Tables
\usepackage{booktabs}
\usepackage{longtable}
\usepackage{array}
\usepackage{multirow}
\usepackage{float}
\usepackage{pdflscape}

% -----------------------------------------------------------------------
% NOTE: Citations are handled by Quarto's pandoc-citeproc engine using
% the CSL file (csl/nlm-citation-sequence-brackets.csl).
% Do NOT add \usepackage{biblatex} here — it conflicts with citeproc.
% -----------------------------------------------------------------------

% Override class caption: use bold label instead of small-caps (sc)
% The class sets labelfont=sc which renders "Figure" as "FIGURE".
\captionsetup{labelfont=bf}

% -----------------------------------------------------------------------
%  Chapter title format: "1. Introduction" on one line (no "Chapter X")
% -----------------------------------------------------------------------
\makeatletter
\renewcommand{\@makechapterhead}[1]{%
  \abovechapterskip
  {\parindent \z@ \chapteralign \normalfont
    \chapterfont
    \ifnum \c@secnumdepth >\m@ne
      \if@mainmatter
        \thechapter.\enspace
      \fi
    \fi
    \interlinepenalty\@M
    #1\par\nobreak
    \chapterbelowskip
  }
  \thispagestyle{\chapter@p@gestyle}
}
\makeatother

% -----------------------------------------------------------------------
%  Suppress Quarto's auto-generated \maketitle — we have a custom title
%  page in tex/frontmatter.tex instead.
% -----------------------------------------------------------------------
\renewcommand{\maketitle}{}

% -----------------------------------------------------------------------
%  Thesis metadata — used by \ttitle, \authorname, \supname, etc.
% -----------------------------------------------------------------------
\thesistitle{Dairy intake and muscle ageing in older women: a prospective
  analysis from the OsteoLaus cohort}
\supervisor{PD Dr.~med.~Dr.~phil.~Patricia Chocano-Bedoya}
\examiner{}
\degree{Master of Science in Bioinformatics and Computational Biology}
% \author is intentionally omitted here; Quarto injects it from _quarto.yml
% which sets \authorname via the class's \author redefinition.
\addresses{}
\subject{Nutritional Epidemiology}
\keywords{sarcopenia, dairy, muscle, ageing, cohort study, OsteoLaus}
\university{University of Fribourg}
\department{Department of Biology}
\group{}
\faculty{Faculty of Science and Medicine}

\AtBeginDocument{
  \hypersetup{pdftitle=\ttitle}
  \hypersetup{pdfauthor=\authorname}
  \hypersetup{pdfkeywords=\keywordnames}
}
\makeatletter
\@ifpackageloaded{bookmark}{}{\usepackage{bookmark}}
\makeatother
\makeatletter
\@ifpackageloaded{caption}{}{\usepackage{caption}}
\AtBeginDocument{%
\ifdefined\contentsname
  \renewcommand*\contentsname{Table of contents}
\else
  \newcommand\contentsname{Table of contents}
\fi
\ifdefined\listfigurename
  \renewcommand*\listfigurename{List of Figures}
\else
  \newcommand\listfigurename{List of Figures}
\fi
\ifdefined\listtablename
  \renewcommand*\listtablename{List of Tables}
\else
  \newcommand\listtablename{List of Tables}
\fi
\ifdefined\figurename
  \renewcommand*\figurename{Figure}
\else
  \newcommand\figurename{Figure}
\fi
\ifdefined\tablename
  \renewcommand*\tablename{Table}
\else
  \newcommand\tablename{Table}
\fi
}
\@ifpackageloaded{float}{}{\usepackage{float}}
\floatstyle{ruled}
\@ifundefined{c@chapter}{\newfloat{codelisting}{h}{lop}}{\newfloat{codelisting}{h}{lop}[chapter]}
\floatname{codelisting}{Listing}
\newcommand*\listoflistings{\listof{codelisting}{List of Listings}}
\makeatother
\makeatletter
\makeatother
\makeatletter
\@ifpackageloaded{caption}{}{\usepackage{caption}}
\@ifpackageloaded{subcaption}{}{\usepackage{subcaption}}
\makeatother
\usepackage{bookmark}
\IfFileExists{xurl.sty}{\usepackage{xurl}}{} % add URL line breaks if available
\urlstyle{same}
\hypersetup{
  pdftitle={Dairy intake and muscle ageing in older women},
  pdfauthor={Anna Boss},
  hidelinks,
  pdfcreator={LaTeX via pandoc}}


\title{Dairy intake and muscle ageing in older women}
\usepackage{etoolbox}
\makeatletter
\providecommand{\subtitle}[1]{% add subtitle to \maketitle
  \apptocmd{\@title}{\par {\large #1 \par}}{}{}
}
\makeatother
\subtitle{A prospective analysis from the OsteoLaus cohort}
\author{Anna Boss}
\date{2026-06-28}
\begin{document}
\maketitle

% -----------------------------------------------------------------------
%  Front matter — inserted before the first Quarto chapter
% -----------------------------------------------------------------------

\frontmatter
\pagestyle{plain}

% -----------------------------------------------------------------------
%  TITLE PAGE
% -----------------------------------------------------------------------
\begin{titlepage}
\begin{center}

\vspace*{.06\textheight}
{\scshape\LARGE \univname\par}\vspace{1.5cm}
\textsc{\Large Master's Thesis}\\[0.5cm]

\HRule \\[0.4cm]
{\huge \bfseries \ttitle\par}\vspace{0.4cm}
\HRule \\[1.5cm]

\begin{minipage}[t]{0.4\textwidth}
\begin{flushleft} \large
\emph{Author:}\\
\authorname
\end{flushleft}
\end{minipage}
\begin{minipage}[t]{0.4\textwidth}
\begin{flushright} \large
\emph{Supervisor:}\\
\supname
\end{flushright}
\end{minipage}\\[3cm]

\vfill
\degreename\\[0.4cm]
\deptname\\[2cm]

\vfill
{\large \today}
\vfill

\end{center}
\end{titlepage}

% -----------------------------------------------------------------------
%  ABSTRACT
% -----------------------------------------------------------------------
\begin{abstract}
\addchaptertocentry{\abstractname}

% Paste your abstract here (≈ 250 words).
% Use \textbf{Background:}, \textbf{Methods:}, \textbf{Results:},
% \textbf{Conclusions:} if a structured abstract is required.

\end{abstract}

% -----------------------------------------------------------------------
%  ACKNOWLEDGEMENTS
% -----------------------------------------------------------------------
\begin{acknowledgements}
\addchaptertocentry{\acknowledgementname}

% Acknowledge supervisors, co-authors, funding, and personal support here.

\end{acknowledgements}

% -----------------------------------------------------------------------
%  TABLE OF CONTENTS / FIGURES / TABLES
% -----------------------------------------------------------------------
\tableofcontents
\listoffigures
\listoftables

% -----------------------------------------------------------------------
%  ABBREVIATIONS
% -----------------------------------------------------------------------
\begin{abbreviations}{ll}

\textbf{ALM}     & \textbf{A}ppendicular \textbf{L}ean \textbf{M}ass \\
\textbf{ALMI}    & \textbf{A}ppendicular \textbf{L}ean \textbf{M}ass \textbf{I}ndex \\
\textbf{BMI}     & \textbf{B}ody \textbf{M}ass \textbf{I}ndex \\
\textbf{CI}      & \textbf{C}onfidence \textbf{I}nterval \\
\textbf{CoLaus}  & \textbf{Co}hort \textbf{Laus}anne \\
\textbf{DXA}     & \textbf{D}ual-energy \textbf{X}-ray \textbf{A}bsorptiometry \\
\textbf{EWGSOP2} & \textbf{E}uropean \textbf{W}orking \textbf{G}roup on
                   \textbf{S}arcopenia in \textbf{O}lder \textbf{P}eople 2 \\
\textbf{HGS}     & \textbf{H}and\textbf{g}rip \textbf{S}trength \\
\textbf{HR}      & \textbf{H}azard \textbf{R}atio \\
\textbf{IQR}     & \textbf{I}nter\textbf{q}uartile \textbf{R}ange \\
\textbf{LMM}     & \textbf{L}inear \textbf{M}ixed \textbf{M}odel \\
\textbf{MET}     & \textbf{M}etabolic \textbf{E}quivalent of \textbf{T}ask \\
\textbf{MPS}     & \textbf{M}uscle \textbf{P}rotein \textbf{S}ynthesis \\
\textbf{OR}      & \textbf{O}dds \textbf{R}atio \\
\textbf{SD}      & \textbf{S}tandard \textbf{D}eviation \\
\textbf{SE}      & \textbf{S}tandard \textbf{E}rror \\

\end{abbreviations}

% -----------------------------------------------------------------------
%  Begin main matter (switches to arabic page numbering)
% -----------------------------------------------------------------------
\mainmatter
\pagestyle{thesis}


\bookmarksetup{startatroot}

\chapter{}\label{section}

\bookmarksetup{startatroot}

\chapter{Introduction}\label{sec-introduction}

\section{The ageing population and the burden of
sarcopenia}\label{the-ageing-population-and-the-burden-of-sarcopenia}

\textbf{The onset of significant skeletal muscle mass decline in adults
represents a complex temporal process that differs between biological
initiation and clinical manifestation.} While sarcopenia often becomes
clinically significant only in the sixth decade of life, the underlying
degenerative processes begin substantially earlier
{[}\citeproc{ref-westerstahl2025}{1},\citeproc{ref-gustafsson2024}{2}{]}.
A large-scale cross-sectional analysis of 50,332 participants aged
45--84 years reveals continuous decline in fat-free thigh muscle volume
across all age groups {[}\citeproc{ref-anderson2025}{3}{]}. This pattern
is also supported by longitudinal evidence
{[}\citeproc{ref-westerstahl2025}{1}{]}.

\textbf{Age-related muscle loss results from multiple interconnected
physiological and metabolic changes, including neural degeneration,
mitochondrial dysfunction, altered amino acid metabolism, and impaired
muscle protein synthesis.} At the neural level, healthy older adults
experience significant motor unit remodelling, with approximately 40\%
fewer motor units by age 71 {[}\citeproc{ref-piasecki2016}{4}{]}. The
progressive loss of muscle mass is fundamentally driven by the dual
mechanisms of muscle fibre atrophy and hypoplasia, largely resulting
from chronic imbalances between muscle protein synthesis (MPS) and
breakdown (MPB) {[}\citeproc{ref-wilkinson2018}{5}{]}.

\section{Postmenopausal women as a particularly vulnerable
population}\label{postmenopausal-women-as-a-particularly-vulnerable-population}

\section{Dairy as a candidate dietary
exposure}\label{dairy-as-a-candidate-dietary-exposure}

\section{Study objective and
hypothesis}\label{study-objective-and-hypothesis}

\bookmarksetup{startatroot}

\chapter{Methods}\label{sec-methods}

\section{Study design and population}\label{study-design-and-population}

I analyzed longitudinal data from the CoLaus\textbar{}PsyCoLaus cohort
and its nested sub-study, OsteoLaus. CoLaus is a population-based cohort
of men and women aged 35 to 75 years at baseline living in Lausanne,
Switzerland, recruited between June 2003 and May 2006 at a single study
center. Follow-up assessments of the CoLaus cohort were conducted at
four intervals, i.e., 2009-2013, 2014-2018, 2018-2021, and 2022-2026.
OsteoLaus is a sub-cohort nested within CoLaus, comprising women who
participated in the CoLaus baseline examination and were invited to an
additional study focused on bone health. The OsteoLaus study included
\(1\,475\) women aged 50 to 80 years at baseline
{[}\citeproc{ref-shevroja2019}{6}{]}.

Eligibility criteria for the OsteoLaus study included written informed
consent, body mass index (BMI) between 15 and 40 kg/m², and availability
of dual-energy X-ray absorptiometry (DXA) data (obtained in vertebrae
and femur) . In addition, ultrasound bone density measurements in the
heel and information on previous osteoporotic fractures had to be
available. Women with known diseases affecting bone metabolism or
skeletal sites were excluded.

Data were collected approximately every 2.5 years at five visits:
baseline (2010-2012, n = \(1\,475\)), second visit (2012-2015, n =
\(1\,349\)), third visit (2015-2018, n = \(1\,242\)), fourth visit
(2017-2020, n = \(1\,104\)), and fifth visit (2020-2022, n = 944). Loss
to follow-up of participants were mainly due to, death,
institutionalisation, refusal to participate, BMI above 40 kg/m², or
severe psychiatric disease {[}\citeproc{ref-vendrami2025a}{7}{]}.

Further details on recruitment and methodology for CoLaus and OsteoLaus
were described by Firmann et al. {[}\citeproc{ref-firmann2008}{8}{]} and
Shevroja et al. {[}\citeproc{ref-shevroja2019}{6}{]}, respectively.
Additional study information is available at the
CoLaus\textbar{}PsyCoLaus
\href{https://www.colaus-psycolaus.ch/}{website}.
\textbf{?@fig-timeline} shows the study population sizes, visit time
intervals, and the variables collected at each visit.

\section{Dataset matching}\label{sec-matching}

Visits from the CoLaus and OsteoLaus cohorts were aligned using a
predefined visit-matching scheme, with OsteoLaus serving as the backbone
dataset. OsteoLaus baseline was paired with CoLaus follow-up 1 (F1),
OsteoLaus visit 3 (V3) with CoLaus follow-up 2 (F2), and OsteoLaus visit
4 (V4) with CoLaus follow-up 3 (F3), corresponding to study time points
T1, T2, and T3, respectively (\textbf{?@fig-timeline}). OsteoLaus visit
2 (V2) was excluded because it had no corresponding CoLaus assessment
and was not part of the longitudinal analysis framework. Similarly, the
baseline visit of CoLaus was excluded. OsteoLaus visit 5 (V5) was
retained as an OsteoLaus-only assessment and defined as time point T4 to
retain two time points for gait speed. For variables available in both
cohorts, the OsteoLaus value was preferentially retained, with the
corresponding CoLaus value used only when the OsteoLaus value was
missing. Variables collected exclusively at CoLaus visits, such as dairy
intake, were retained unchanged. The final dataset included the time
difference between the matched CoLaus and OsteoLaus examination dates
for each study time point.

\section{Exposure}\label{exposure}

\subsection{Dietary assessment}\label{dietary-assessment}

Dietary intake, including dairy consumption, was assessed in the CoLaus
study from the first follow-up visit onwards using a self-administered
semi-quantitative food frequency questionnaire (FFQ). The FFQ was
developed and validated in Geneva and includes 97 food items covering
intake over the previous 4 weeks
{[}\citeproc{ref-bernstein1995}{9},\citeproc{ref-marques-vidal2018}{10}{]}.
Participants reported consumption frequency and portion size for each
item. Portion sizes were defined as seen in \textbf{?@tbl-ffq} for dairy
products. Daily intake (g/day or mL/day) was calculated by combining
frequency and portion size information. Energy and nutrient values were
derived using the French CIQUAL food composition database
{[}\citeproc{ref-ciqual}{11}{]}.

\begin{longtable}[]{@{}
  >{\raggedright\arraybackslash}p{(\linewidth - 6\tabcolsep) * \real{0.1644}}
  >{\raggedright\arraybackslash}p{(\linewidth - 6\tabcolsep) * \real{0.3973}}
  >{\raggedright\arraybackslash}p{(\linewidth - 6\tabcolsep) * \real{0.0822}}
  >{\raggedright\arraybackslash}p{(\linewidth - 6\tabcolsep) * \real{0.3562}}@{}}
\caption{Food Frequency Questionnaire (FFQ) items and portion sizes for
women.}\tabularnewline
\toprule\noalign{}
\begin{minipage}[b]{\linewidth}\raggedright
FFQ number
\end{minipage} & \begin{minipage}[b]{\linewidth}\raggedright
Food item
\end{minipage} & \begin{minipage}[b]{\linewidth}\raggedright
Unit
\end{minipage} & \begin{minipage}[b]{\linewidth}\raggedright
Portion size (S / E / L)
\end{minipage} \\
\midrule\noalign{}
\endfirsthead
\toprule\noalign{}
\begin{minipage}[b]{\linewidth}\raggedright
FFQ number
\end{minipage} & \begin{minipage}[b]{\linewidth}\raggedright
Food item
\end{minipage} & \begin{minipage}[b]{\linewidth}\raggedright
Unit
\end{minipage} & \begin{minipage}[b]{\linewidth}\raggedright
Portion size (S / E / L)
\end{minipage} \\
\midrule\noalign{}
\endhead
\bottomrule\noalign{}
\endlastfoot
1 & Plain yogurt & g & 90 / 180 / 270 \\
2 & Low-fat yogurt & g & 90 / 180 / 270 \\
3 & Fruit/aroma yogurt & g & 90 / 180 / 270 \\
4 & Cottage cheese (0\% fat) & g & 60 / 125 / 200 \\
5 & Cottage cheese, ricotta & g & 50 / 100 / 150 \\
6 & Feta, mozzarella & g & 20 / 40 / 60 \\
7 & Gruyère, tomme, camembert & g & 30 / 60 / 90 \\
8 & Cheese fondue & g & 175 / 350 / 525 \\
52 & Butter & g & 4 / 7 / 10 \\
53 & Cream (35\% fat) & g & 8 / 15 / 23 \\
63 & Cream tart, cream cake & g & 50 / 100 / 150 \\
68 & Ice cream, sorbet & g & 50 / 100 / 150 \\
71 & Butter (for cooking) & g & 4 / 7 / 10 \\
82 & Milk in coffee (0\% fat) & mL & 15 / 30 / 45 \\
83 & Milk in coffee (non 0\% fat) & mL & 15 / 30 / 45 \\
84 & Coffee creamer & mL & 6 / 12 / 18 \\
85 & Milk as drink (0\% fat) & mL & 75 / 150 / 225 \\
86 & Milk as drink (non 0\% fat) & mL & 75 / 150 / 225 \\
\end{longtable}

FFQ = food frequency questionnaire; S = small, E = medium, L = large
portion. Portion sizes provided by .

\subsection{Total dairy intake}\label{total-dairy-intake}

Total dairy intake was defined as the sum of milk, cheese, yogurt,
cream, butter, and dairy-based desserts consumed per day. Intake was
expressed in g/day, assuming a density of 1 g/mL for liquid dairy
products to allow

zation across units. Ice cream (FFQ68) was reported together with
sorbet. The entire amount of FFQ68 reported was assumed to represent a
dairy-based product and was included in total dairy intake.

\subsection{Dairy subtypes}\label{dairy-subtypes}

Dairy subtypes were derived by summing predefined FFQ items. Fermented
dairy products included yogurt and cheese. Non-fermented dairy products
included milk, cream, and ice cream. Full-fat dairy products comprised
full-fat milk, yogurt, cheese, cream, butter, and ice cream, while
low-fat dairy products included skimmed milk, low-fat yogurt, and
reduced-fat cheese. Detailed FFQ item mappings are provided in the
Supplementary Material.

\subsection{Adherence to dietary
guidelines}\label{adherence-to-dietary-guidelines}

Adherence to Swiss Society for Nutrition recommendations for dairy
intake was classified as meeting or not meeting recommended intake
levels {[}\citeproc{ref-dairyguid}{12}{]}. Intake was calculated using
weighted frequency scores based on portion size (0.5 for smaller, 1.0
for equal, 1.5 for larger portions) for selected dairy items
contributing substantially to total intake (FFQ1 to FFQ8, FFQ85 and
FFQ86).

\subsection{Exposure modelling}\label{exposure-modelling}

Total dairy intake was the primary exposure and was modeled both as a
continuous variable and in quartiles. Quartile cut-offs were derived
from the distribution at the first follow-up visit and applied to all
subsequent visits. Secondary analyses explored dairy subtypes, modeled
exclusively as continuous variables.

For longitudinal outcomes (handgrip strength and ALMI), I used the
cumulative average of total dairy intake, calculated as the mean of all
available previous intakes up to each visit. For gait speed, I utilized
a lagged cumulative average (mean of previous intakes, excluding the
most recent visit) to account for the absence of dairy data at the T4
time point. For sarcopenia incidence, the Cox proportional hazards model
incorporated both continuous and quartile-based intake measures.
Sensitivity analyses were conducted to test the robustness of our
findings, including the assessment of adherence to dietary guidelines
and the use of quartiles in mixed linear models.

\section{Covariates}\label{covariates}

\subsection{Covariate selection}\label{covariate-selection}

Covariates were selected based on prior evidence of direct or indirect
associations with muscle mass, muscle strength, physical function,
sarcopenia, and dairy intake. A directed acyclic graph (DAG) was used to
represent the assumed relationships between variables (Supplementary
Figure XX ). The DAG guided the selection of a minimal sufficient
adjustment set to control for confounding. Variables that did not
contribute to blocking confounding paths, as well as potential
mediators, colliders, or redundant variables, were not included in the
final models.

\subsection{Dietary assessment}\label{dietary-assessment-1}

Total energy intake was assessed as a covariate to account for overall
dietary quantity, since individuals with higher energy intake tend to
consume more of all food groups, including dairy products. Adjusting for
total energy intake helps separate the association of dairy intake from
overall food consumption and reduces confounding due to differences in
dietary energy intake between participants. Total energy intake was
calculated as the sum of all FFQ items (all 97 food items), using energy
values derived from the French CIQUAL food composition table
{[}\citeproc{ref-ciqual}{11}{]}.

\subsection{Sociodemographic factors}\label{sociodemographic-factors}

Sociodemographic factors are associated with dietary patterns and
lifestyle. Therefore, these factors are also associated with health
outcome. The sociodemographic variables were self-reported. I included
the following variables:

\begin{itemize}
\item
  Age of each participant was derived at each visit using the birth date
  and examination date.
\item
  Educational level was defined as the highest level of education
  completed at the time of recruitment. It was grouped into three
  categories based on the International Standard Classification of
  Education (ISCED): low (ISCED 0-2), medium (ISCED 3), and high (ISCED
  4-8) {[}\citeproc{ref-ISCED}{13}{]}.
\item
  Marital status was recorded during the CoLaus visits and included the
  category `living alone' and `living in a partnership'.
\end{itemize}

The following sociodemographic variables were not included as
confounding factor:

\begin{itemize}
\item
  Sex and geographical location were not used since the entire cohort
  consisted of females living in Lausanne, Switzerland.
\item
  Ethnicity was not included since most participants (\textgreater{}
  98\%) were Caucasian.
\item
  Income and occupation were not included because the data were not
  available.
\end{itemize}

\subsection{Anthropometric factors}\label{anthropometric-factors}

Higher BMI is often linked to higher energy intake and to lower muscle
function. Thus, BMI was considered as covariate. Height and weight were
measured using standardized protocols with a portable stadiometer (Seca
216) and electronic scale (Seca Clara 803), respectively. Measurements
were taken with participants barefoot and in light clothing, with
precision of 0.1 cm and 0.1 kg, respectively. BMI was calculated as
weight divided by height squared (kg/m²) and categorized according to
WHO definitions: underweight (\textless{} 18.5 kg/m²), normal weight
(18.5--24.9 kg/m²), overweight (25.0--29.9 kg/m²), and obese (\(\geq\)
30.0 kg/m²).

\subsection{Lifestyle}\label{lifestyle}

Smoking may indirectly affect muscle mass through lower levels of
physical activity {[}\citeproc{ref-degens2015}{14}{]}. It may also
directly contribute to the loss of muscle mass by reducing protein
synthesis and increasing muscle breakdown through metabolic pathway
changes. Thus, the smoking status of the participant was included as
covariate. Smoking history of participants was collected through
self-report. Participants were grouped into three categories:

\begin{itemize}
\tightlist
\item
  Current smokers: those who were actively smoking at the time of the
  study, regardless of how much they smoked
\item
  Former smokers: those who had smoked in the past but had quit before
  the study, regardless of how long ago they had quit
\item
  Never smokers: those who had never smoked
\end{itemize}

Smoking status was recorded at each CoLaus study visit. Some reported
changes were not realistic, for example, switching from never smoker to
former or current smoker, or from former or current smoker to never
smoker. In these cases, the ``never smoker'' status was changed to
``former smoker'', assuming that participants had answered based on
their current situation rather than their lifetime history.

Alcohol consumption is not consistently linked to lower grip strength or
appendicular lean mass. However, some evidence suggests that higher
intake of alcohol may be associated with lower muscle mass in a
dose-dependent manner
{[}\citeproc{ref-skinner2023}{15},\citeproc{ref-cui2019}{16}{]}. Alcohol
comsumtion was therefore included as a covariate.

Alcohol consumption was assessed using two self-reported questionnaires.
The method used in the main analysis asked participants how many
alcoholic drinks they had consumed in the previous 7 days measured in
units/week and was used for the main analysis. The second method asked
about the number of glasses of wine, bottles of beer, or shots of
spirits consumed in the past month measured in gram ethanol/day and was
used for sensitivity analyses.

Based on these responses, the number of units consumed per week was
calculated. According to Swiss national guidelines, one standard drink
was defined as one glass of wine, one bottle of beer, or one shot of
spirits, each containing about 10 g to 12 g of pure ethanol
{[}\citeproc{ref-standardglas}{17}{]}. The alcohol intake in g/day was
converted into units per week using 10 g of ethanol per unit.
Participants were then grouped into four categories based on their total
weekly alcohol intake:

\begin{itemize}
\tightlist
\item
  Non-drinkers: 0 units/week (0 g/day), which included both lifetime and
  current abstainers.
\item
  Light drinkers: \(\geq\) 1 to 3 units/week (\textgreater{} 0--6
  g/day), representing occasional consumption.
\item
  Moderate drinkers: \(\geq\) 4 to 7 units/week (\textgreater{} 6--12
  g/day), representing frequent consumption within the Swiss low-risk
  limit for women.
\item
  Heavy drinkers: \textgreater{} 7 units/week (\textgreater{} 12 g/day),
  categorized as hazardous consumption exceeding the Swiss national
  guidelines for women.
\end{itemize}

Physical activity (PA) was considered as a covariate because it strongly
correlates with muscle function outcomes in older adults. Systematic
reviews and longitudinal studies consistently show that higher levels of
physical activity are associated with better muscle strength, power, and
functional performance in older populations
{[}\citeproc{ref-ramsey2021}{18},\citeproc{ref-ward-ritacco2022}{19}{]}.

PA was measured using a Physical Activity Frequency Questionnaire (PAFQ)
developed in 1998 for the Geneva population and containing 70 types of
physical activity {[}\citeproc{ref-bernstein1998}{20}{]}. The
questionnaire recorded how many days during the past week (0-7 days)
each activity was performed and the duration per day (0-10 h in 15 min
increments) {[}\citeproc{ref-verhoog2019}{21}{]}. Each activity was
classified according to its metabolic equivalent of task (MET) value as
sedentary behaviour (\textless{} 2 METs), light physical activity (2 to
\textless{} 3 METs), moderate physical activity (3-6 METs), or vigorous
physical activity (\textgreater{} 6 METs).

This assessment was only available at the first and second follow-up
visits of the CoLaus study. A previous study in a subpopulation of
CoLaus showed moderate correlation between physical activity measured
with the PAFQ and accelerometer-based measurements
{[}\citeproc{ref-verhoog2019}{21}{]}. Self-reported questionnaires tend
to overestimate physical activity and underestimate sedentary behaviour.
Therefore, physical activity was analyzed as tertiles representing low,
moderate, and high activity levels rather than as a continuous variable.
Tertiles were used because no established reference values were
available.

The tertiles were defined using the combined amount of moderate and
vigorous physical activity reported at the first follow-up visit and
were subsequently applied to all time points. Only moderate and vigorous
physical activity levels were used because these measures showed the
smallest difference compared with accelerometer-based assessments
{[}\citeproc{ref-verhoog2019}{21}{]}.

\subsection{Clinical variables}\label{clinical-variables}

Type 2 diabetes was assessed as a covariate because previous studies
have shown a bidirectional relationship between diabetes and sarcopenia.
Diabetes may contribute to reduced muscle mass and function through
insulin resistance, chronic inflammation, oxidative stress, and altered
protein metabolism, while low muscle mass may also worsen glucose
metabolism and increase diabetes risk {[}\citeproc{ref-chen2023}{22}{]}.

Diabetes status was assessed using three sources recorded during CoLaus
visits: self-report, fasting plasma glucose (FPG), and glycated
haemoglobin (HbA1c). Self-report was based on the question ``Have you
ever been told that you had diabetes?''. Among the objective measures,
HbA1c was considered the primary indicator due to its higher reliability
and was used whenever available. If HbA1c was missing, FPG was used. If
neither biomarker was available, diabetes status was based on
self-report. Description of the sample collection procedure can be found
elsewhere {[}\citeproc{ref-ducraux2021}{23}{]}.

Hormonal replacement therapy (HRT) status was assessed for a potential
covariate because previous research shows mixed but biologically
plausible effects on muscle strength, mass, and function in
postmenopausal women {[}\citeproc{ref-xu2020}{24}{]}. Given the
variability in HRT use among participants, controlling for this factor
helps ensure that the observed associations are not influenced by
hormonal differences.

HRT use was assessed using self-reported questionnaire data.
Participants were asked whether they were currently using HRT and at
what age HRT was started. Based on these responses, HRT status was
classified into three categories: current HRT use, past HRT use, and
never or not current HRT use. Participants who answered ``Yes'' to
current HRT use were classified as current users. Participants who
answered ``No'' but provided a valid age at HRT initiation were
classified as past users. Participants who answered ``No'' and had no
indication of previous HRT use were classified as never or not current
users. HRT status was set to missing when no information on current HRT
use was available.

Hypertension was considered as a covariant because of the established
association between hypertension and reduced muscle mass
{[}\citeproc{ref-han2020}{25}{]}. Hypertension status was derived using
a hierarchical approach based on three sources of information: measured
hypertension, documented antihypertensive treatment, and self-reported
antihypertensive medication use. Measured hypertension was defined as
systolic blood pressure \(\geq\) 140 mmHg and/or diastolic blood
pressure \(\geq\) 90 mmHg and was considered the primary source.
Participants were classified as having hypertension if hypertension was
measured. If measured hypertension data were missing, documented
antihypertensive treatment was used. If both measured hypertension and
documented treatment data were unavailable, self-reported
antihypertensive medication use was considered. Participants were
classified as not having hypertension if at least one data source was
available and none indicated hypertension. Hypertension status was set
to missing if all three sources were unavailable.

Cardiovascular disease (CVD) is associated with reduced muscle mass and
function via multiple mechanisms, including reduced physical activity
due to symptoms such as fatigue and dyspnoea, impaired muscle perfusion
and oxygen delivery, systemic inflammation, and shared metabolic risk
factors such as insulin resistance and age-related pathways.

CVD event status was derived as a composite binary variable (any CVD
event: yes/no) based on 13 self-reported and clinically reported
cardiovascular conditions, including myocardial infarction, stroke,
heart failure, coronary artery disease, angina, cardiomyopathy, cardiac
surgery, heart valve surgery, treated arrhythmia, peripheral vascular
surgery, cardiac catheterisation, coronary artery bypass graft, and
percutaneous coronary intervention.

Participants were classified as having a CVD event if at least one of
the component conditions was reported as ``Yes''. Participants were
classified as having no CVD event if at least one component was
available, and all reported conditions were ``No''. Data were classified
as missing if no information was available across all component
variables. CVD status was available only at baseline.

\subsection{Medication and supplement
use}\label{medication-and-supplement-use}

The use of prescribed medication and supplements was assessed by
self-report during CoLaus visits. Each drug type was labelled in line
with the Anatomical Therapeutic Chemical (ATC) classification. Some
drugs have beneficial effects on muscle mass and muscle function,
whereas others negatively impact muscle-related outcomes.

Lipid-lowering medications have been linked to muscle-related adverse
effects in clinical reports {[}\citeproc{ref-shin2026}{26}{]}. However,
evidence for an association with reduced muscle mass or impaired muscle
function remains inconsistent, with some studies reporting negative
associations {[}\citeproc{ref-shin2026}{26}{]} and others finding no
effect {[}\citeproc{ref-fazel2026}{27}{]}. Lipid-lowering drugs were
defined as all medications classified under ATC code C10.

Use of systemic corticosteroids and benzodiazepines negatively impact
muscle strength, mass, and functional outcomes in older adults. Thus,
ATC codes starting with H02 were included for use of systemic
corticosteroids and N05B for use of benzodiazepines.

The evidence for muscle-protective effects of bisphosphonates in humans
is limited, but it was suggested that bone-to-muscle paracrine
signalling may mediate benefits {[}\citeproc{ref-klein2018}{28}{]}.
Therefore, all participants reporting ACT code M05BA were considered to
take bisphosphonates.

Supplementation with vitamin D (ACT code A11) and calcium (ACT code
A12A) is positively associated with muscle outcomes
{[}\citeproc{ref-gao2026}{29}{]}. Thus, these variables were included as
covariates in the analysis.

\section{Outcome Variables}\label{outcome-variables}

\subsection{Handgrip strength}\label{handgrip-strength}

Handgrip strength (HGS) was measured at all CoLaus visits using a JAMAR
Baseline® hydraulic hand dynamometer (Fabrication Enterprises, Inc.,
White Plains, NY, USA) following the guidelines of the American Society
of Hand Therapists. Values were recorded in pounds and converted to
kilograms using a conversion factor of 0.454.

In OsteoLaus, HGS was assessed only at visit V5 using the same protocol.
The highest value of three measurements of the dominant hand was used
for analysis. Measurements were recorded directly in kilograms. Further
details of the procedure are described in Vendrami et al.~(2025)
{[}\citeproc{ref-vendrami2025}{30}{]}.

\subsection{Appendicular lean mass
(ALM)}\label{appendicular-lean-mass-alm}

ALM was defined as the sum of lean muscle mass of the arms and legs. ALM
was measured at all five OsteoLaus visits using dual-energy X-ray
absorptiometry (DXA). Two devices were used over time: the Hologic
Horizon A System™ for the first two visits and the GE Lunar iDXA™ system
for the remaining visits.

A study within the same cohort revealed that the Hologic Horizon A
System generally leads to a higher lean muscle mass (1293.7
g \(\pm\) 1643.5 g) {[}\citeproc{ref-vendrami2024}{31}{]}. However,
measured lean mass of the arms and legs was higher for the Lunar iDXA
device (arms 151.4 g \(\pm\) 274.8 g; legs: 777.1 g \(\pm\) 490.0 g).

Because systematic differences were observed between devices, ALM values
were harmonized using previously published cross-calibration equations
derived within the same OsteoLaus cohort
{[}\citeproc{ref-vendrami2024}{31}{]}. Hologic-derived values were
converted to Lunar-equivalent values prior to analysis using the
following equation: \(ALM_{Holigic}=−369.3+0.967\cdot ALM_{Lunar}\).

Appendicular lean mass index (ALMI) was calculated as ALM divided by
height squared (kg/m²). ALM was converted from grams to kilograms and
height from centimetres to metres before calculation. Sensitivity
analyses additionally included alternative muscle mass index based on
ALM adjusted for body mass index (ALM/BMI).

\subsection{Gait speed}\label{gait-speed}

Gait speed was used as a measure of physical performance and functional
capacity. Gait speed was assessed using a 6-metre walk test performed at
usual walking pace and expressed in meters per second (m/s). Detailed
procedures are described in Vendrami et al.
{[}\citeproc{ref-vendrami2025}{30}{]}. Slow gait speed was defined as
\(≤\) 0.8 m/s.

\subsection{Sarcopenia}\label{sarcopenia}

Sarcopenia was defined according to the criteria of the European Working
Group on Sarcopenia in Older People 2 (EWGSOP2)
{[}\citeproc{ref-cruz-jentoft2019}{32}{]} and was used for the main
analyses. Muscle strength was defined as low when handgrip strength was
\textless~16 kg, muscle mass as low when ALMI was \textless~5.5 kg/m²,
and physical performance as low when gait speed was \(≤\)~0.8~m/s. These
values represent the female-specific thresholds of EWGSOP2 sarcopenia
definition {[}\citeproc{ref-cruz-jentoft2019}{32}{]}.

Participants were classified into four categories: no sarcopenia,
probable sarcopenia (low muscle strength only), confirmed sarcopenia
(low muscle strength and low muscle mass), and severe sarcopenia (low
muscle strength, low muscle mass, and low physical performance). Because
gait speed was measured only at visits V4 and V5, severe sarcopenia
could not be assessed at earlier visits. This limitation did not affect
the primary sarcopenia classification, as participants were classified
as sarcopenic if they met the criteria for either confirmed or severe
sarcopenia.

In sensitivity analyses, sarcopenia was additionally defined using the
Foundation for the National Institutes of Health (FNIH) criteria
{[}\citeproc{ref-studenski2014}{33}{]}, based on the combination of low
muscle strength (HGS \textless~16 kg) and low muscle mass relative to
body mass index (ALM/BMI \textless~0.512). Participants were classified
as having sarcopenia or no sarcopenia in line with this definition.

\section{Statistical Methods}\label{statistical-methods}

All analyses were performed in R version 4.5.1
{[}\citeproc{ref-Rprogram}{34}{]}. Data preparation and manipulation was
done with dplyr (version 1.2.0) {[}\citeproc{ref-wickham2026}{35}{]}.
Results were visualized using ggplot2 (version 4.0.2)
{[}\citeproc{ref-wickham2016a}{\textbf{wickham2016a?}}{]} For
descriptive analyses, participants were grouped into quartiles based on
baseline daily total dairy intake (g/day). Categorical variables are
presented as counts and percentages, and continuous variables as mean ±
standard deviation (SD) or median with interquartile range (IQR), as
appropriate.

\subsection{Missing data}\label{missing-data}

Missing data were handled using multiple imputation by chained equations
(MICE) implemented in the R package mice (version 3.19.0)
{[}\citeproc{ref-buuren2011}{36}{]}. Imputation was performed under a
fully conditional specification framework and with the assumption that
data were missing at random.

Longitudinal data were reshaped into wide format (one row per
participant with separate columns for each visit), allowing
within-person correlations across time to be incorporated in the
imputation model. A total of 10 imputed datasets were generated, with 30
iterations per dataset.

Continuous variables were imputed using predictive mean matching. Binary
variables were imputed using logistic regression, and ordinal variables
using proportional odds models. Derived variables were not imputed
directly. Instead, they were recomputed after imputation from the
underlying variables. This approach preserves internal consistency
between raw and derived variables. Auxiliary variables (e.g.~participant
identifiers, examination dates, and additional clinical or lifestyle
variables) were included as predictors to improve imputation accuracy
but were not imputed as such.

For OsteoLaus, visit structure and measurement design were explicitly
incorporated. Variables were only imputed at visits where measurement
was possible (e.g.~gait speed only at T3 and T4). For DXA-derived
variables, measurements obtained from different devices were prevented
from directly predicting each other to reduce systematic bias across
measurement methods.

Predictor matrices were manually specified to preserve temporal
structure. Only missing values were imputed. Observed values were left
unchanged. Post-imputation diagnostics included trace plots, density
plots, and strip plots to assess convergence and distributional
plausibility across imputations.

\subsection{Data cleaning and
exclusions}\label{data-cleaning-and-exclusions}

\subsection{\texorpdfstring{Covariate selection
}{Covariate selection }}\label{covariate-selection-1}

\subsection{Linear mixed-effects
models}\label{linear-mixed-effects-models}

\subsection{\texorpdfstring{\textbf{Cox proportional hazards
model}}{Cox proportional hazards model}}\label{cox-proportional-hazards-model}

\section{Ethics}\label{ethics}

CoLaus \textbar{} PsyCoLaus and the OsteoLaus studies were approved by
the Institutional Ethics Committee of the University of Lausanne,
Switzerland. All participants signed a written informed consent form
after having received a detailed description of the goal and funding of
the studies.

\hfill\break

\section{Data and Code Availability
Statement}\label{data-and-code-availability-statement}

This observational cohort study is reported in accordance with the
STROBE guidelines {[}\citeproc{ref-STROBE}{37}{]}.

The data used in this study are not publicly available but may be
requested subject to approval. More information on data request can be
found in \url{www.colaus-psycolaus.ch}.

All code used for data processing and statistical analyses is available
on GitHub: \url{https://github.com/ann4boss/OsteoLaus-Dairy-Sarcopenia}.

\bookmarksetup{startatroot}

\chapter{Results}\label{sec-results}

\section{Study population}\label{study-population}

\section{Baseline characteristics}\label{sec-baseline}

\section{Dairy intake and handgrip strength}\label{sec-hgs}

\section{Dairy intake and appendicular lean mass index}\label{sec-almi}

\section{Dairy intake and gait speed}\label{sec-gait}

\section{Incident sarcopenia}\label{sec-cox}

\section{Sensitivity analyses}\label{sec-sensitivity}

\bookmarksetup{startatroot}

\chapter{Discussion}\label{sec-discussion}

\section{Principal findings}\label{principal-findings}

\section{Comparison with existing
literature}\label{comparison-with-existing-literature}

\section{Biological mechanisms}\label{biological-mechanisms}

\section{Strengths}\label{sec-strengths}

\section{Limitations}\label{sec-limitations}

\section{Implications for public health}\label{sec-implications}

\bookmarksetup{startatroot}

\chapter{Conclusion}\label{sec-conclusion}

\bookmarksetup{startatroot}

\chapter*{References}\label{references}
\addcontentsline{toc}{chapter}{References}

\markboth{References}{References}

\cleardoublepage
\phantomsection
\addcontentsline{toc}{part}{Appendices}
\appendix

\chapter*{Supplementary Material}\label{sec-supplementary}
\addcontentsline{toc}{chapter}{Supplementary Material}

\markboth{Supplementary Material}{Supplementary Material}

\section*{Directed Acyclic Graph (DAG)}\label{sec-dag}
\addcontentsline{toc}{section}{Directed Acyclic Graph (DAG)}

\markright{Directed Acyclic Graph (DAG)}

\section*{Consort flow diagram}\label{sec-consort}
\addcontentsline{toc}{section}{Consort flow diagram}

\markright{Consort flow diagram}

\begin{figure}

\centering{

\includegraphics[width=0.9\linewidth,height=\textheight,keepaspectratio]{appendices/../images/consort_flow.png}

}

\caption{\label{fig-consort}Consort flow diagram}

\end{figure}%

\section*{Supplementary tables}\label{sec-sup-tables}
\addcontentsline{toc}{section}{Supplementary tables}

\markright{Supplementary tables}

\section*{Supplementary analyses}\label{sec-sup-analyses}
\addcontentsline{toc}{section}{Supplementary analyses}

\markright{Supplementary analyses}

\phantomsection\label{refs}
\begin{CSLReferences}{0}{1}
\bibitem[\citeproctext]{ref-westerstahl2025}
\CSLLeftMargin{1. }%
\CSLRightInline{Maria Westerståhl, Gustav Jörnåker, Eva Jansson, Ulrika
Aasa, Michael Ingre, Kaveh Pourhamidi, Brun Ulfhake, Thomas Gustafsson.
Rise and {Fall} of {Physical Capacity} in a {General Population}: {A}
47-{Year Longitudinal Study}. Journal of Cachexia, Sarcopenia and
Muscle. 2025;16(6):e70134.
doi:\href{https://doi.org/10.1002/jcsm.70134}{10.1002/jcsm.70134}}

\bibitem[\citeproctext]{ref-gustafsson2024}
\CSLLeftMargin{2. }%
\CSLRightInline{Thomas Gustafsson, Brun Ulfhake. Aging {Skeletal
Muscles}: {What Are} the {Mechanisms} of {Age-Related Loss} of
{Strength} and {Muscle Mass}, and {Can We Impede Its Development} and
{Progression}? International Journal of Molecular Sciences. 2024
Oct;25(20).
doi:\href{https://doi.org/10.3390/ijms252010932}{10.3390/ijms252010932}}

\bibitem[\citeproctext]{ref-anderson2025}
\CSLLeftMargin{3. }%
\CSLRightInline{David B. Anderson, Aaron J. Beach, Lingxiao Chen, Henry
J. Feng, Marnee J. McKay, Zachary A. Smith, Kenneth A. Weber, Evert Onno
Wesselink, James M. Elliott. What is normal age-related thigh muscle
composition among 45- to 84-year-old adults from the {UK Biobank} study.
GeroScience. 2025 Feb;47(1):1175--85.
doi:\href{https://doi.org/10.1007/s11357-024-01304-y}{10.1007/s11357-024-01304-y}}

\bibitem[\citeproctext]{ref-piasecki2016}
\CSLLeftMargin{4. }%
\CSLRightInline{Mathew Piasecki, Alex Ireland, David A. Jones, Jamie S.
McPhee. Age-dependent motor unit remodelling in human limb muscles.
Biogerontology. 2016 Jun;17(3):485--96.
doi:\href{https://doi.org/10.1007/s10522-015-9627-3}{10.1007/s10522-015-9627-3}}

\bibitem[\citeproctext]{ref-wilkinson2018}
\CSLLeftMargin{5. }%
\CSLRightInline{D. J. Wilkinson, M. Piasecki, P. J. Atherton. The
age-related loss of skeletal muscle mass and function: {Measurement} and
physiology of muscle fibre atrophy and muscle fibre loss in humans.
Ageing Research Reviews. 2018 Nov;47:123--32.
doi:\href{https://doi.org/10.1016/j.arr.2018.07.005}{10.1016/j.arr.2018.07.005}}

\bibitem[\citeproctext]{ref-shevroja2019}
\CSLLeftMargin{6. }%
\CSLRightInline{Enisa Shevroja, Pedro Marques-Vidal, Bérengère
Aubry-Rozier, Gabriel Hans, Fernando Rivadeneira, Olivier Lamy, Didier
Hans. Cohort {Profile}: {The OsteoLaus} study. International Journal of
Epidemiology. 2019 Aug;48(4):1046--1047g.
doi:\href{https://doi.org/10.1093/ije/dyy276}{10.1093/ije/dyy276}}

\bibitem[\citeproctext]{ref-vendrami2025a}
\CSLLeftMargin{7. }%
\CSLRightInline{Colin Vendrami, Elena Gonzalez Rodriguez, Guillaume
Gatineau, Peter Vollenweider, Pedro Marques-Vidal, Olivier Lamy, Didier
Hans, Enisa Shevroja. Prediction of 10-{Year Fragility Fractures Using
Muscle Health Indicators} in {Postmenopausal Women}: {The OsteoLaus
Cohort}. Journal of Cachexia, Sarcopenia and Muscle. 2025
Jun;16(3):e13837.
doi:\href{https://doi.org/10.1002/jcsm.13837}{10.1002/jcsm.13837} PubMed
PMID: \href{https://www.ncbi.nlm.nih.gov/pubmed/40464199}{40464199};
PubMed Central PMCID:
\href{https://www.ncbi.nlm.nih.gov/pmc/articles/PMC12134769}{PMC12134769}.}

\bibitem[\citeproctext]{ref-firmann2008}
\CSLLeftMargin{8. }%
\CSLRightInline{Mathieu Firmann, Vladimir Mayor, Pedro Marques Vidal,
Murielle Bochud, Alain Pécoud, Daniel Hayoz, Fred Paccaud, Martin
Preisig, Kijoung S. Song, Xin Yuan, Theodore M. Danoff, Heide A.
Stirnadel, Dawn Waterworth, Vincent Mooser, Gérard Waeber, Peter
Vollenweider. The {CoLaus} study: A population-based study to
investigate the epidemiology and genetic determinants of cardiovascular
risk factors and metabolic syndrome. BMC Cardiovascular Disorders. 2008
Mar;8(1):6.
doi:\href{https://doi.org/10.1186/1471-2261-8-6}{10.1186/1471-2261-8-6}}

\bibitem[\citeproctext]{ref-bernstein1995}
\CSLLeftMargin{9. }%
\CSLRightInline{Bernstein, Huot, Morabia. Amélioration des performances
d'un questionnaire alimentaire semi-quantitatif comparé à un rappel des
24 heures. Sante Publique. 1995.}

\bibitem[\citeproctext]{ref-marques-vidal2018}
\CSLLeftMargin{10. }%
\CSLRightInline{Pedro Marques-Vidal, Gérard Waeber, Peter Vollenweider,
Idris Guessous. Socio-demographic and lifestyle determinants of dietary
patterns in {French-speaking Switzerland}, 2009--2012. BMC Public
Health. 2018 Jan;18(1):131.
doi:\href{https://doi.org/10.1186/s12889-018-5045-1}{10.1186/s12889-018-5045-1}}

\bibitem[\citeproctext]{ref-ciqual}
\CSLLeftMargin{11. }%
\CSLRightInline{French Agency for Food, Environmental and Occupational
Health \& Safety. Ciqual. French food composition table.
https://ciqual.anses.fr/.}

\bibitem[\citeproctext]{ref-dairyguid}
\CSLLeftMargin{12. }%
\CSLRightInline{Federal Food Safety and Veterinary Office FSVO. Swiss
{Dietary Recommendations}.
https://www.blv.admin.ch/blv/en/home/lebensmittel-und-ernaehrung/ernaehrung/empfehlungen-informationen/schweizer-ernaehrungsempfehlungen.html.}

\bibitem[\citeproctext]{ref-ISCED}
\CSLLeftMargin{13. }%
\CSLRightInline{International {Standard Classification} of {Education}
({ISCED}).
https://ec.europa.eu/eurostat/statistics-explained/index.php?title=International\_Standard\_Classification\_of\_Education\_(ISCED).}

\bibitem[\citeproctext]{ref-degens2015}
\CSLLeftMargin{14. }%
\CSLRightInline{Hans Degens, Ghislaine Gayan-Ramirez, Hieronymus W. H.
van Hees. Smoking-induced {Skeletal Muscle Dysfunction}. {From Evidence}
to {Mechanisms}. American Journal of Respiratory and Critical Care
Medicine. 2015 Mar;191(6):620--5.
doi:\href{https://doi.org/10.1164/rccm.201410-1830PP}{10.1164/rccm.201410-1830PP}}

\bibitem[\citeproctext]{ref-skinner2023}
\CSLLeftMargin{15. }%
\CSLRightInline{Alcohol {Consumption} and {Measures} of {Sarcopenic
Muscle Risk}: {Cross-Sectional} and {Prospective Associations Within}
the {UK Biobank Study} \textbar{} {Calcified Tissue International}
\textbar{} {Springer Nature Link}.
https://link.springer.com/article/10.1007/s00223-023-01081-4.}

\bibitem[\citeproctext]{ref-cui2019}
\CSLLeftMargin{16. }%
\CSLRightInline{Yufei Cui, Cong Huang, Haruki Momma, Shota Sugiyama,
Kaijun Niu, Ryoichi Nagatomi. The longitudinal association between
alcohol consumption and muscle strength: {A} population-based
prospective study. Journal of Musculoskeletal \& Neuronal Interactions.
2019;19(3):294--9. PubMed PMID:
\href{https://www.ncbi.nlm.nih.gov/pubmed/31475936}{31475936}; PubMed
Central PMCID:
\href{https://www.ncbi.nlm.nih.gov/pmc/articles/PMC6737549}{PMC6737549}.}

\bibitem[\citeproctext]{ref-standardglas}
\CSLLeftMargin{17. }%
\CSLRightInline{Das {Standardglas Alkohol}... - {Alcohol Facts}.
https://www.aktionstag-alkoholprobleme.ch/de/das-standardglas-alkohol.}

\bibitem[\citeproctext]{ref-ramsey2021}
\CSLLeftMargin{18. }%
\CSLRightInline{Keenan A. Ramsey, Anna G. M. Rojer, Luke D'Andrea, René
H. J. Otten, Martijn W. Heymans, Marijke C. Trappenburg, Sjors Verlaan,
Anna C. Whittaker, Carel G. M. Meskers, Andrea B. Maier. The association
of objectively measured physical activity and sedentary behavior with
skeletal muscle strength and muscle power in older adults: {A}
systematic review and meta-analysis. Ageing Research Reviews. 2021
May;67:101266.
doi:\href{https://doi.org/10.1016/j.arr.2021.101266}{10.1016/j.arr.2021.101266}}

\bibitem[\citeproctext]{ref-ward-ritacco2022}
\CSLLeftMargin{19. }%
\CSLRightInline{Christie L. Ward-Ritacco, Mary L. Greaney, Phillip G.
Clark, Deborah Riebe. The {Longitudinal Association Between Physical
Activity} and {Physical Function} in {Older Women}. Frontiers in Sports
and Active Living. 2022 Jul;4.
doi:\href{https://doi.org/10.3389/fspor.2022.879025}{10.3389/fspor.2022.879025}}

\bibitem[\citeproctext]{ref-bernstein1998}
\CSLLeftMargin{20. }%
\CSLRightInline{M. Bernstein, D. Sloutskis, S. Kumanyika, A. Sparti, Y.
Schutz, A. Morabia. Data-based {Approach} for {Developing} a {Physical
Activity Frequency Questionnaire}. American Journal of Epidemiology.
1998 Jan;147(2):147--54.
doi:\href{https://doi.org/10.1093/oxfordjournals.aje.a009427}{10.1093/oxfordjournals.aje.a009427}}

\bibitem[\citeproctext]{ref-verhoog2019}
\CSLLeftMargin{21. }%
\CSLRightInline{Sanne Verhoog, Cédric Gubelmann, Idris Guessous, Arjola
Bano, Oscar H. Franco, Pedro Marques-Vidal. Comparison of the {Physical
Activity Frequency Questionnaire} ({PAFQ}) with accelerometry in a
middle-aged and elderly population: {The CoLaus} study. Maturitas. 2019
Nov;129:68--75.
doi:\href{https://doi.org/10.1016/j.maturitas.2019.08.004}{10.1016/j.maturitas.2019.08.004}
PubMed PMID:
\href{https://www.ncbi.nlm.nih.gov/pubmed/31547917}{31547917}.}

\bibitem[\citeproctext]{ref-chen2023}
\CSLLeftMargin{22. }%
\CSLRightInline{Huiling Chen, Xiaojing Huang, Meiyuan Dong, Song Wen,
Ligang Zhou, Xinlu Yuan. The {Association Between Sarcopenia} and
{Diabetes}: {From Pathophysiology Mechanism} to {Therapeutic Strategy}.
Diabetes, Metabolic Syndrome and Obesity. 2023 May;16:1541--54.
doi:\href{https://doi.org/10.2147/DMSO.S410834}{10.2147/DMSO.S410834}}

\bibitem[\citeproctext]{ref-ducraux2021}
\CSLLeftMargin{23. }%
\CSLRightInline{Pauline Ducraux, Gérard Waeber, Pedro Marques-Vidal. Do
{Weight} trajectories influence diabetes control? {A} prospective study
in {Switzerland} ({CoLaus} study). Preventive Medicine Reports. 2021
Jun;23:101473.
doi:\href{https://doi.org/10.1016/j.pmedr.2021.101473}{10.1016/j.pmedr.2021.101473}
PubMed PMID:
\href{https://www.ncbi.nlm.nih.gov/pubmed/34258179}{34258179}; PubMed
Central PMCID:
\href{https://www.ncbi.nlm.nih.gov/pmc/articles/PMC8259406}{PMC8259406}.}

\bibitem[\citeproctext]{ref-xu2020}
\CSLLeftMargin{24. }%
\CSLRightInline{Yang Xu, Kai-Li Deng, Tian-Fang Xing, Ya-Qing Mei,
Su-Mei Xiao. Effect of hormone therapy on muscle strength in
postmenopausal women: A systematic review and meta-analysis of
randomized controlled trials. Menopause. 2020 Jul;27(7):827.
doi:\href{https://doi.org/10.1097/GME.0000000000001538}{10.1097/GME.0000000000001538}}

\bibitem[\citeproctext]{ref-han2020}
\CSLLeftMargin{25. }%
\CSLRightInline{Moon K. Han, Raeda Anderson, Emilie Viennois, Didier
Merlin. Examination of food consumption in {United States} adults and
the prevalence of inflammatory bowel disease using {National Health
Interview Survey} 2015. PloS one. 2020;15(4):e0232157.
doi:\href{https://doi.org/10.1371/journal.pone.0232157}{10.1371/journal.pone.0232157}
PubMed PMID:
\href{https://www.ncbi.nlm.nih.gov/pubmed/32324818}{32324818}; PubMed
Central PMCID:
\href{https://www.ncbi.nlm.nih.gov/pmc/articles/PMC7179926}{PMC7179926}.}

\bibitem[\citeproctext]{ref-shin2026}
\CSLLeftMargin{26. }%
\CSLRightInline{Ah-Reum Shin, Keun-Gyu Park, Sung-Woo Kim. Association
{Between Lipid-Lowering Drug Use} and {Sarcopenia}: {Analysis} of the
{Korea National Health} and {Nutrition Examination Survey}. Journal of
Lipid and Atherosclerosis. 2026 Jan;15(1):151--60.
doi:\href{https://doi.org/10.12997/jla.2026.15.1.151}{10.12997/jla.2026.15.1.151}}

\bibitem[\citeproctext]{ref-fazel2026}
\CSLLeftMargin{27. }%
\CSLRightInline{Seyedeh Dalia Fazel, Lisanne Tap, Paolo Fabbietti,
Andrea Corsonello, Robert Kob, Cornel C. Sieber, Regina
Roller-Wirnsberger, Gerhard Wirnsberger, Ronit Ben-Romano, Itshak
Melzer, Tomasz Kostka, Agnieszka Guligowska, Axel C. Carlsson, Francesc
Formiga, Rafael Moreno-González, Johan Ärnlöv, Fabrizia Lattanzio,
Francesco Mattace-Raso, other SCOPE investigators. Statin use does not
impair muscle health in older adults: Findings from the {SCOPE} study.
Aging Clinical and Experimental Research. 2026 Feb;38(1):89.
doi:\href{https://doi.org/10.1007/s40520-026-03337-z}{10.1007/s40520-026-03337-z}}

\bibitem[\citeproctext]{ref-klein2018}
\CSLLeftMargin{28. }%
\CSLRightInline{Gordon L. Klein. The {Role} of {Bone Secreted Factors}
in {Burn-Induced Muscle Cachexia}. Current Osteoporosis Reports. 2018
Feb;16(1):26--31.
doi:\href{https://doi.org/10.1007/s11914-018-0416-1}{10.1007/s11914-018-0416-1}}

\bibitem[\citeproctext]{ref-gao2026}
\CSLLeftMargin{29. }%
\CSLRightInline{Xiao-Xue Gao, Hong-Wei Li, Dan Liu, Yu-Wei He, Wei Ji,
Ting-Ting Liang, Jiu-Wei Cui. Vitamin {D} and calcium regulate body
muscle function and fat accumulation: {A} retrospective study based on
the {UK} biobank. Journal of Nutritional Oncology. 2026 Mar;11(1):10.
doi:\href{https://doi.org/10.1097/JN9.0000000000000063}{10.1097/JN9.0000000000000063}}

\bibitem[\citeproctext]{ref-vendrami2025}
\CSLLeftMargin{30. }%
\CSLRightInline{Colin Vendrami, Elena Gonzalez Rodriguez, Guillaume
Gatineau, Peter Vollenweider, Pedro Marques-Vidal, Olivier Lamy, Didier
Hans, Enisa Shevroja. Prevalence and incidence of sarcopenia in {Swiss}
postmenopausal women: Findings from the {OsteoLaus Cohort}. Swiss
Medical Weekly. 2025 Jan;155(1):4034--4.
doi:\href{https://doi.org/10.57187/s.4034}{10.57187/s.4034}}

\bibitem[\citeproctext]{ref-vendrami2024}
\CSLLeftMargin{31. }%
\CSLRightInline{Colin Vendrami, Guillaume Gatineau, Elena Gonzalez
Rodriguez, Olivier Lamy, Didier Hans, Enisa Shevroja. Standardization of
body composition parameters between {GE Lunar iDXA} and {Hologic Horizon
A} and their clinical impact. JBMR Plus. 2024 Sep;8(9):ziae088.
doi:\href{https://doi.org/10.1093/jbmrpl/ziae088}{10.1093/jbmrpl/ziae088}}

\bibitem[\citeproctext]{ref-cruz-jentoft2019}
\CSLLeftMargin{32. }%
\CSLRightInline{Alfonso J Cruz-Jentoft, Gülistan Bahat, Jürgen Bauer,
Yves Boirie, Olivier Bruyère, Tommy Cederholm, Cyrus Cooper, Francesco
Landi, Yves Rolland, Avan Aihie Sayer, Stéphane M Schneider, Cornel C
Sieber, Eva Topinkova, Maurits Vandewoude, Marjolein Visser, Mauro
Zamboni. Sarcopenia: Revised {European} consensus on definition and
diagnosis. Age and Ageing. 2019 Jan;48(1):16--31.
doi:\href{https://doi.org/10.1093/ageing/afy169}{10.1093/ageing/afy169}
PubMed PMID:
\href{https://www.ncbi.nlm.nih.gov/pubmed/30312372}{30312372}; PubMed
Central PMCID:
\href{https://www.ncbi.nlm.nih.gov/pmc/articles/PMC6322506}{PMC6322506}.}

\bibitem[\citeproctext]{ref-studenski2014}
\CSLLeftMargin{33. }%
\CSLRightInline{Stephanie A. Studenski, Katherine W. Peters, Dawn E.
Alley, Peggy M. Cawthon, Robert R. McLean, Tamara B. Harris, Luigi
Ferrucci, Jack M. Guralnik, Maren S. Fragala, Anne M. Kenny, Douglas P.
Kiel, Stephen B. Kritchevsky, Michelle D. Shardell, Thuy-Tien L. Dam,
Maria T. Vassileva. The {FNIH Sarcopenia Project}: {Rationale}, {Study
Description}, {Conference Recommendations}, and {Final Estimates}. The
Journals of Gerontology Series A: Biological Sciences and Medical
Sciences. 2014 May;69(5):547--58.
doi:\href{https://doi.org/10.1093/gerona/glu010}{10.1093/gerona/glu010}
PubMed PMID:
\href{https://www.ncbi.nlm.nih.gov/pubmed/24737557}{24737557}; PubMed
Central PMCID:
\href{https://www.ncbi.nlm.nih.gov/pmc/articles/PMC3991146}{PMC3991146}.}

\bibitem[\citeproctext]{ref-Rprogram}
\CSLLeftMargin{34. }%
\CSLRightInline{R Core Team. R: A language and environment for
statistical computing {[}Manual{]}. Vienna, Austria: R Foundation for
Statistical Computing; 2025.}

\bibitem[\citeproctext]{ref-wickham2026}
\CSLLeftMargin{35. }%
\CSLRightInline{Hadley Wickham, Romain François, Lionel Henry, Kirill
Müller, Davis Vaughan. Dplyr: A grammar of data manipulation. 2026.}

\bibitem[\citeproctext]{ref-buuren2011}
\CSLLeftMargin{36. }%
\CSLRightInline{Stef van Buuren, Karin Groothuis-Oudshoorn. Mice:
{Multivariate Imputation} by {Chained Equations} in {R}. Journal of
Statistical Software. 2011 Dec;45:1--67.
doi:\href{https://doi.org/10.18637/jss.v045.i03}{10.18637/jss.v045.i03}}

\bibitem[\citeproctext]{ref-STROBE}
\CSLLeftMargin{37. }%
\CSLRightInline{Strengthening the {Reporting} of {Observational Studies}
in {Epidemiology} ({STROBE}): {Explanation} and {Elaboration}. Annals of
Internal Medicine. 2009 Feb.}

\end{CSLReferences}




\end{document}
